\documentclass[a4paper,12pt]{article}

\geometry{left=20mm,right=20mm,top=25mm,bottom=20mm} % задание полей текста
\usepackage{wrapfig}
%% Стиль колонтитулов
% \fancyhead[RO,LE]{\hyperlink{intro}{Содержание}} % Right odd,  Left even
\fancyhead[LO]{\@lecture}        % Right even, Left odd
\fancyhead[R]{}

\fancyfoot[RO,LE]{\thepage}         % Right odd,  Left even
\fancyfoot[RE,LO]{\CourseName}      % Right even, Left odd
\fancyfoot[C]{}
% Un~comment these to erase foot (and comment footrulewidth renewcommand)
%\fancyfoot{}
%\fancyhead[C]{-~\thepage~-}
\renewcommand{\footrulewidth}{0.4pt}
\usepackage{textcomp}

% Новая команда \lecture{№ лекции}{название}
% После этой команды весь текст до следующей такой же команды будет
% принадлежать конкретной лекции, имя которой будет в колонтитуле каждой страницы
\usepackage{xifthen}
\def\@lecture{}%
\newcommand{\lecture}[2]{
    \ifthenelse{\isempty{#2}}{%
        \def\@lecture{Лекция #1}%
    }{%
        \def\@lecture{Лекция #1: #2}%
    }%
    \section{\@lecture}
}

\def\@lecture{}%
\newcommand{\question}[2]{
    \ifthenelse{\isempty{#2}}{%
        \def\@lecture{Билет #1}%
    }{%
        \def\@lecture{Билет #1: #2}%
    }%
    \section{\@lecture}
}
% ------------ Text settings ------------
%%% Гиппер ссылки
\renewcommand{\linkcolor}{blue}
\renewcommand{\citecolor}{green}
\renewcommand{\filecolor}{magenta}
\renewcommand{\urlcolor}{NavyBlue}

\usepackage{multicol}	   % Для текста в нескольких колонках

% -----------  Images -----------
\graphicspath{{images/}{img/}{figures/}{fig/}}  % Путь к папкам с картинками
\newcommand{\figL}[3]{%      Для быстрой вставки картинок
    \begin{figure}[h!]
        \centering
        \includegraphics[width=#2\textwidth]{#1}
        \label{fig:#3}
    \end{figure}%
}
\newcommand{\fig}[2]{%    
    \begin{figure}[h!]
        \centering
        \includegraphics[width=#2\textwidth]{#1}
    \end{figure}%
}

% ----------- Math short-cats
\newcommand{\R}{\ensuremath{\mathbb{R}}}
\newcommand{\N}{\ensuremath{\mathbb{N}}}
\newcommand{\Cx}{\ensuremath{\mathbb{C}}}
\newcommand{\Z}{\ensuremath{\mathbb{Z}}}
\newcommand{\E}{\ensuremath{\mathbb{E}}}
\newcommand{\Q}{\ensuremath{\mathbb{Q}}}
\def\CB{\mathcal{B}}
\def\CC{\mathcal{C}}
\def\CE{\mathcal{E}}
\def\CR{\mathcal{R}}
\def\CA{\mathcal{A}}
% \def\CF{\mathcal{F}}
\def\CG{\mathcal{G}}
\def\CS{\mathcal{S}}
\def\CD{\mathcal{D}}
\def\CH{\mathcal{H}}
\def\CP{\mathcal{P}}
\def\CM{\mathcal{M}}
\def\FB{\mathfrak{B}}

% You can write your commands below
\usepackage{gensymb}
\usepackage{enumitem}
\usepackage{amsmath}
\newcommand{\F}{\ensuremath{\mathcal{F} }}
\newcommand{\Anu}{\ensuremath{\mathcal{A}_{\nu}}}
\DeclareMathOperator{\FDU}{FDU}

% ----------- Math and theorems -----------
\usepackage[many]{tcolorbox}
\usepackage{mdframed}
\usepackage[dvipsnames]{xcolor}

\newtheorem*{remark}      {Замечание}
\newtheorem*{next0}      {Следствие}
\newtheorem*{next1}      {Следствие 1}
\newtheorem*{next2}      {Следствие 2}
\theoremstyle{definition}
\newtheorem{lemma}{Лемма}[section]
\newtheorem{claim}[lemma]{Утверждение}
\newtheorem{theorem_}[lemma]{Теорема}
\newenvironment{theorem}%
{\begin{mdframed}[backgroundcolor=black!30!white!30]
        % \setlength{\topsep}{-\parskip}\setlength{\partopsep}{0pt}
        \begin{theorem_}}%
            {\end{theorem_}\end{mdframed}}

\theoremstyle{definition}
\newtheorem{definition}[lemma]{Определение}
\tcolorboxenvironment{definition}{
    enhanced,
    borderline={0.8pt}{0pt}{gray!70},
    borderline={0.4pt}{2pt}{black},
    boxrule=0.4pt,
    colback=pink!70!white!30,
    coltitle=black,
    sharp corners
}
\usepackage{ulem}
\newcommand{\mdef}[1]
{\! \textit{\uwave{\textcolor{red!10!black}{#1}}}}
% http://dkhramov.dp.ua/Comp.LatexCyrillicFonts#.XMrWLegzaUk


%https://tex.stackexchange.com/questions/223694/how-to-draw-a-text-box-with-shadow-borders-in-latex

\newtheorem{exercise_}[lemma]{Пример}
\newenvironment{exercise}%
{\begin{tcolorbox}[enhanced,width=\textwidth,center upper,drop fuzzy shadow southwest,
            colframe=red!50!black,colback=orange!100!yellow!50!white!30]
        \begin{exercise_}}%
            {\end{exercise_}\end{tcolorbox}}

\usepackage{dashrule}
\renewenvironment{proof}{\smallskip{\noindent\Large
        \color{red!50!black}\itshape$\lhd$ \normalsize Начало доказательства
        \hdashrule[0.5ex]{0.65\textwidth}{0.5mm}{3mm 3pt 1mm 2pt} }
    \small\itshape}{{\color{red!50!black}
            \hdashrule[0.5ex]{0.66\textwidth}{0.5mm}{3mm 3pt 1mm 2pt}
            \normalsize Конец доказательства \Large$\rhd$}}

\newenvironment{solution}{\smallskip{\noindent\Large
        \color{red!50!black}\itshape$\lhd$ \normalsize Начало решения
        \hdashrule[0.5ex]{0.75\textwidth}{0.5mm}{3mm 3pt 1mm 2pt} }
    \small}{{\color{red!50!black}
            \hdashrule[0.5ex]{0.75\textwidth}{0.5mm}{3mm 3pt 1mm 2pt}
            \normalsize\itshape Конец решения \Large$\rhd$}}

% Обводка кружочком множеств
\usepackage{tikz}
\usetikzlibrary{decorations.pathreplacing,shapes.misc,patterns}
\newcommand*\circled[1]{\tikz[baseline=(char.base)]{
        \node[shape=circle,draw,inner sep=2pt] (char) {#1};}}

\counterwithin*{equation}{section}

\usepackage{stackrel}
\newsavebox\MBox
\newcommand\Cline[2][red]{{\sbox\MBox{$#2$}%
            \rlap{\usebox\MBox}\color{#1}\rule[-1.2\dp\MBox]{\wd\MBox}{0.5pt}}}

%%% Всю шаблонную информацию можно менять тут
\newcommand{\FullCourseNameFirstPart}{\so{ДИСКРЕТНЫЙ АНАЛИЗ}}
\newcommand{\FullCourseNameSecondPart}{\so{}}
\newcommand{\SemesterNumber}{III}
\newcommand{\SchoolName}{ФПМИ}
\newcommand{\TrackName}{ПМИ}
\newcommand{\LecturerInitials}{Райгородский Андрей Михайлович}
\newcommand{\AutherInitials}{Герасимов Фёдор}
\newcommand{\VKLink}{https://vk.com/gefedya}
\newcommand{\OverleafLink}{https://www.overleaf.com/read/qfpvwgrfrfkr}
\newcommand{\GithubLink}{https://github.com/MIPT-Group/Lectures_Tex_Club}

\newcommand{\CourseName}{Краткое Название} %  Используется в преамбуле для создания названия предмета в верхнем контитуле   
\newcommand{\CourseDate}{Осень 2020}           %  Используется в преамбуле для создания даты в нижнем контитуле и в title_page

%\includeonly{lectures/lect05,lectures/lect06}  % Чтобы скомпилировать только часть лекций

\begin{document}
    %%% Всю шаблонную информацию можно менять тут
\newcommand{\FullCourseNameFirstPart}{\so{АРХИТЕКТУРА КОМПЬЮТЕРОВ И}}
\newcommand{\FullCourseNameSecondPart}{\so{ОПЕРАЦИОННЫЕ СИСТЕМЫ}}
\newcommand{\SchoolName}{ФПМИ}
\newcommand{\TrackName}{ПМИ}
\newcommand{\SemesterNumber}{III}
\newcommand{\LecturerInitials}{Яковлев Виктор Вадимович}
\newcommand{\AutherInitials}{Савелий Романов}
\newcommand{\VKLink}{https://vk.com/romanov_savelij}
\newcommand{\GithubLink}{https://github.com/MIPT-Group/Lectures_Tex_Club}

\begin{titlepage}
	\clearpage\thispagestyle{empty}
	\centering
	
	\textit{Федеральное государственное автономное учреждение \\
		высшего профессионального образования}
	\vspace{0.5ex}
	
	\textbf{Московский Физико-Технический Институт \\ КЛУБ ТЕХА ЛЕКЦИЙ}
	\vspace{20ex}
	
	\rule{\linewidth}{0.5mm}
	{\textbf{\FullCourseNameFirstPart}}
	\\
	{\textbf{\FullCourseNameSecondPart}}
	\rule{\linewidth}{0.5mm}
	
	\SemesterNumber\ СЕМЕСТР
	\\
	Физтех Школа: \textit{\SchoolName}
	\\
	Направление: \textit{\TrackName}
	\\
	Лектор: \textit{\LecturerInitials}
	\vspace{1ex}
	
	\begin{figure}[!ht]
		\centering
		\includegraphics[width=0.6\textwidth, height=0.3\textheight]{logo_LTC}
		\label{fig:my_label}
	\end{figure}
\begin{flushright}
	\noindent
	Автор: \href{\VKLink}{\textit{\AutherInitials}}
	\\
	\href{\GithubLink}{\textit{Проект на github}} % Опционально, если хотите учавствовать в рейтинге
\end{flushright}
	
	\vfill
	\CourseDate\ года
	\pagebreak
	
\end{titlepage}
    \newpage
    \hypertarget{intro}{}
    \tableofcontents
    \newpage
    
    \section{3 Семестр}
    \subsection{Оценки комбинаторных величин}

\subsubsection{Асимптотика биномиальных коэффициентов}

\begin{Def}
	Функции $f, g:N \rightarrow \mathbb{R}$ называются \emph{aсимптотически равными}, если \(\frac{f(n)}{g(n)} \rightarrow 1, n \rightarrow \infty \). Обозначение \(f(n) \sim g(n)\).
\end{Def}


\begin{example}
	В прошлом семестре из простых соображений бинома Ньютона была доказана оценка:
	\[ \frac{2^{2n}}{2n + 1} < C^n_{2n} < 2^{2n} \]
	
	Уже эта оценка позволяет нам получить асимптотику логарифма биномиального коэффициента:
	
	\[ 2n\ln(n) - \ln(2n+1) < \ln(C^n_{2n}) < 2n\ln(n) \]
	
	
	\[2 n \ln 2\left(1-\frac{\ln (2 n+1)}{2 n \ln 2}\right) < \ln C_{2 n}^{n} < 2n \ln 2, \]
	
	
	откуда \( \ln(C^n_{2n}) \sim 2n\ln(2) \).
	Из налогичных рассуждений следует оценка \(C^n_{2n} \sim \left(4 + o(1)\right)^n\).\footnote{Запись $f\left( n\right) = \left( g\left( n\right) + o\left( 1\right) \right) ^n $ эквивалентна $\sqrt[n]{f(n)} \sim g\left( n\right) $}
\end{example}

%\begin{note}
%	Пусть есть функция $f: \N \rightarrow \R$. Тогда запись $f\left( n\right) = \left( g\left( n\right) + o\left( 1\right) \right) ^n $ эквивалентна $\sqrt[n]{f(n)} \sim g\left( n\right) $.
%\end{note}

\begin{theorem}[Формула Стирлинга (б/д)]
	\[ n! \sim \sqrt{2\pi n} (\frac{n}{e})^n \]
\end{theorem}

\begin{note}
	Из \emph{Формулы Стирлинга} следует оценка для того же биномиального коэффициента: \(C^n_{2n} \sim \frac{4^n}{\sqrt{\pi n}}\).
\end{note}

\begin{theorem}
	Для $a \in \left( 0, 1\right)$ верно:
	\[C_{n}^{[a n]}=\left(a^{-a}(1-a)^{a-1}+o(1)\right)^{n} \]
	
\end{theorem}

\begin{proof}
	Распишем:
	
	%\begin{align*}
		%C_{n}^{[a n]}=\frac{n!}{[a n]!(n-[a n])!} &=
		%\frac{\sqrt{2 \pi n}\left(\frac{n}{e}\right)^{n}}{\sqrt{2 \pi[a n]}\left(\frac{[a n]}{e}\right)^{[a n]} \sqrt{2 \pi(n-[a n])}\left(\frac{(n-[a n])}{e}\right)^{(n-[a n])}} =
	%\end{align*}
	
	$$C_{n}^{[a n]}=\frac{n!}{[a n]!(n-[a n])!} =
	\frac{\sqrt{2 \pi n}\left(\frac{n}{e}\right)^{n}}{\sqrt{2 \pi[a n]}\left(\frac{[a n]}{e}\right)^{[a n]} \sqrt{2 \pi(n-[a n])}\left(\frac{(n-[a n])}{e}\right)^{n-[a n]}}= $$
	$$
	= \frac{\sqrt{2 \pi n}}{\sqrt{2 \pi[a n]} \sqrt{2 \pi\left( n-[a n]\right) }} \cdot \left( \frac{[an]}{n}\right) ^{-[an]}\left( \frac{n-[an]}{n}\right) ^{[an]-n}$$
	%=\left(\frac{1}{a^{a}(1-a)^{1-a}}\right)^{n} P(n),$$
	
	\qquad
	
	Тогда
	
	$$
	\sqrt[n]{C_{n}^{[a n]}}=\sqrt[n]{\frac{\sqrt{2 \pi n}}{\sqrt{2 \pi[a n]} \sqrt{2 \pi\left( n-[a n]\right) }}} \cdot \left( \frac{[an]}{n}\right) ^{-\frac{[an]}{n}}\left( \frac{n-[an]}{n}\right) ^{\frac{[an]-n}{n},}
	% \ln \left(\frac{1}{(a)^{a}(1-a)^{1-a}}\right) \sim n \ln \left(\frac{1}{(a)^{a}(1-a)^{1-a}}\right)
	$$
	
	При $n \to \infty$ первый множитель, очевидно, имеет пределом 1, а у второго и третьего и основания, и показатели степени имеют конечный предел. Из свойств экспонетны и логарифма заключаем треубемое. (Если говорить строго, то для экспоненциальной функции имеем $\lim_{x \to a} \exp(x) = \exp(a)$. Переходя к логарифму, для $a > 0$ имеем $\lim_{x \to a} \ln{x} = \ln(a)$. Тогда для $a > 0, b > 0$ имеем $\lim_{n \to \infty} f(n)^{g(n)} = \lim_{n \to \infty} \exp(\ln(f(n))g(n)) = \exp\left( \lim_{n \to \infty}(\ln(f(n))g(n))\right)  = \exp(\ln(\lim_{n \to \infty} f(n))\lim_{n \to \infty} g(n))$).
	
\end{proof}

\begin{theorem}
	Верны следующие оценки:
	\begin{enumerate}
		\item $C_{n}^{k} \leqslant \frac{n^{k}}{k !}$
		\item $C_{n}^{k} = \frac{n^{k}}{k!} \cdot e^{-\frac{k(k-1)}{2n} + O\left(\frac{k^3}{n^2}\right)}$
		\item Если $k = o(n^{\frac{2}{3}})$, то $C_{n}^{k} \sim \frac{n^{k}}{k!} \cdot e^{-\frac{k(k-1)}{2n}}$
		\item Если $k = o(\sqrt{n})$, то $C_{n}^{k} \sim \frac{n^{k}}{k!}$
	\end{enumerate}
	
\end{theorem} \label{exponent}

\begin{proof}
	\begin{enumerate}
		
		\item $C_{n}^{k}=\frac{n(n-1) \ldots(n-k+1)}{k !} \leqslant \frac{n^{k}}{k !} $
		\item Используя разложение функции $e^x$ в ряд Тейлора, получаем:
		
		\begin{eqnarray*}
		C_{n}^{k}=\frac{n(n-1) \ldots(n-k+1)}{k !}=\frac{n^{k}}{k!} \cdot e^{\:\ln \frac{n(n-1) \ldots(n-k+1)}{n^{k}}} \leqslant
		\frac{n^{k}}{k !} \cdot e^{\:ln(1 - \frac{1}{n}) + ln(\frac{2}{n})+\ldots+ln(\frac{k-1}{n})} = \\
		=\frac{n^{k}}{k!} \cdot e^{-\frac{1}{n} + O\left( \frac{1}{n^2}\right)  - \ldots  - \frac{k-1}{n} + O\left( \frac{( k-1)^2}{n^2}\right) } = \frac{n^{k}}{k!} \cdot e^{-\frac{k(k-1)}{2n} + O\left( \frac{k^3}{n^2}\right) }
		\end{eqnarray*}
		
	\end{enumerate}
\end{proof}

\begin{brainer}
	Найти асимптотику выражения:  $ \sum_{i = 0}^{n} \left( C_{n}^{i}\right)^\alpha $
\end{brainer}


\subsubsection{Графы}

\begin{Def}
	Графом называется пара множеств $(V, E)$, где $V$ -- множество каких-то
	объектов, а $E \subset \binom{V}{2}$ -- множество пар объектов из $V$.
\end{Def} 


\begin{Def}
	Маршрутом в графе $G=(V, E)$ называется последовательность $v_{1} e_{1} v_{2} \ldots e_{n} v_{n+1}$, где $e_i = \left\lbrace v_{n}, v_{i+1} \right\rbrace $.
	
	\begin{itemize}
		\item Если $v_{1}=v_{n+1}$, то маршрут называется \emph{замкнутым}.
		Если все $e_{i}$ в замкнутом маршруте различны, то он называется \emph{циклом}, а незамкнутый - \emph{путём}.
		\item Цепь (цикл) называется простой(-ым), если все $v_i$ в нем различны.
	\end{itemize}
\end{Def} 

\begin{Def}
	Граф называется \emph{связным}, если любые две его вершины соединены маршрутом.
\end{Def}

\begin{Def}
	\emph{Деревом} называется связный граф без циклов.
\end{Def}

\begin{Def}
	Вершина степени 1 в дереве называется \emph{листом} или \emph{висячей вершиной}.
\end{Def}

\begin{note}
	В каждом дереве на $n>1$ вершинах есть висячая вершина. (А на самом деле их хотя бы две!)
\end{note}

\begin{theorem} \label{equivalences}
	Следующие свойства графа эквивалентны:
	\begin{enumerate}
		\item Граф является деревом
		\item Граф, у которого любые две вершины соединены единственным маршрутом
		\item Связный граф, у которого число ребер на единицу меньше числа вершин
		\item Ациклический граф, у которого число ребер на единицу меньше числа вершин
	\end{enumerate}
\end{theorem}

\begin{proof}
	\;
	
	1 $\Rightarrow$ 2. 
	Докажем транспозицию.
	Если между двумя вершинами в графе $G$ нету пути, то он не является связным. Если путь между двумя вершинами не единственный, то, пройдя по первому, а затем по второму в обратном  направлении, мы получим цикл.
	
	2 $\Rightarrow $ 3.
	Очевидно, что $G$ связен. Докажем, что $|E|=|V| - 1$ по индукции по числу вершин. База $|V|=1$ очевидна. Предположим, что $|V|=n$, а утверждение верно для всех $k < n$. Удалим из
	графа $G$ некоторое ребро. Т.к. между любыми двумя вершинами существует ровно один простой путь, то $G$ распался на 2 компоненты связности. Применим предположение индукции для каждой из них и получим требуемое.
	
	3 $\Rightarrow $ 4.
	Предположим, что в графе, удовлетворяющем условию 3, есть цикл. Удалим в цикле прозвольное ребро. Граф остался связным, а количество ребер в нем стало равно $n-2$. Но тогда граф несвязный, поскольку в связном графе Будем повторять операцию, пока в графе не останется циклов. Тогда полученный граф будет удовлетворять условию 2, значит и условию 3. Но количество ребер в нем по предположению меньше $n - 1$. Следовательно, наше изначальное предположение было неверным, и граф ациклический.
	
	4 $\Rightarrow $ 1.
	Предположим противное. Будем действовать аналогично предыдущему пункту, удаляя по ребру в цикле на каждом шаге. Граф при этом остается связным. Когда граф станет ациклическим, он станет деревом, но с количеством вершин, меньшим, чем $|V| - 1$. Поскольку $1 \Rightarrow 4$, приходим к противоречию.
\end{proof}

\begin{theorem}[Кэли] \label{Caley}
	Количество деревьев на n занумерованных вершинах равно $n^{n - 2}$.
\end{theorem}

Приведем строгое доказательство теоремы.
\begin{proof}
	Для $n = 1$ утверждение теоремы тривиально.
	Условимся доказывать теорему для $n \geqslant 2$.
	\medskip
	
	Пусть дано дерево $T$ на $n$ вершинах. Возьмем копию дерева $T$ и будем применять следующий алгоритм, пока количество вершин в копии $n^{\prime} \geqslant 3$:
	\begin{enumerate}
		\item Выберем висячую вершину $v$ с минимальным номером. 
		\item Добавим в код Прюфера номер вершины, смежной с $v$.
		\item Удалим из дерева вершину $v$ и выходящее из нее ребро.
		\item Применим алгоритм к модифицированному графу на $n^{\prime} - 1$ вершинах.
	\end{enumerate}
	
	Докажем корректность алгоритма:
	\begin{enumerate} \renewcommand{\theenumi}{\bfseries \Roman{enumi}}
		\item Докажем, что на каждом шаге алгоритм определен. В начале алгоритма граф является деревом, следовательно, в нем есть висячая вершина, а значит и вершина с минимальным номером. Осталось показать, что после применения шага алогритма граф остается деревом. Действительно, при удалении висячей вершины граф остается связным, а цикла возникнуть не может.
		
		\item Докажем, что по завершении алгоритма мы получаем код из $n-2$ символа из элементов $\{1, \ldots, n\}$. Поскольку в начале алгоритма у графа $n$ вершин, а останавливается он при $n = 2$, то количество шагов алгоритма будет равно $n - 2$.
	\end{enumerate}
	
	
	Итак, мы определили соответсвие между $Tree(n)$ и $Prufer(n)$. Оно является отображением. Действительно, если одному дереву соответсвуют разные коды Прюфера, то возьмем первый символ, в котором они различны. Но на этом шаге мы записали в код номер вершины, смежной с висячей с наименьшим номером, которая, очевидно, одна.
	\medskip
	
	Рассмотрим полученное отображение $\varphi : Tree(n) \mapsto Prufer(n)$. Докажем по индукции, что $\varphi$ инъективно:
	
	База $|V| = 2$ очевидна. Предположим, что два различных дерева $T_{1}, T_{2}$ отвечают одному коду $v_{1} \ldots v_{n}, v_{i} \in\{1, \ldots, n\}$.
	Возможны следующие случаи:
	
	\begin{enumerate}
		\item Листы с наименьшим номером в $T_{1}$ и $T_{2}$ различны.
		ННО $k$ - наименьший номер висячей вершины  в $T_1$. Тогда первый код для $T_1$ номер этой вершины не войдет, а в код для $T_2$ - войдет, поскольку она не является висячей в $T_2$. Следовательно, у двух деревьев разные коды Прюфера.
		\item Листы с наименьшим номером совпадают, но различны их соседи. Тогда их коды отличаются уже в первом символе.
		\item Если же совпадают листы с наименьшими номерами и их соседи, то первый символ кодов Прюфера $T_1$ и $T_2$ совпадают, но после вычеркивания остаются два дерева, коды которых различны по предположению
		индукции.
	\end{enumerate}
	
	Докажем теперь, что отображение сюръективно или, что то же самое, предъявим алгоритм восстановления дерева по коду Прюфера.
	Вначале мы имеем последователньость $ \tau = v_{1} \ldots v_{n-2}, v_{i} \in\{1, \ldots, n\}$. размера $n-2$ и множество всех вершин графа: $ V = \left\lbrace 1… n\right\rbrace $. Применим следующюю процедуру:
	
	\begin{enumerate} 
		\item Берем первый элемент кода Прюфера $\tau$, и в множестве $V$ ищем вершину с наименьшим номером, которая не содержится в коде. Найденная вершина и текущий элемент массива с кодом Прюфера составляют ребро дерева. Запишем пару элементов в множество вершин будущего графа $E$.
		\item Удалим номер вершины из множества $V$ и первый элемент коде Прюфера.
		\item Применим алгоритм к модифицированному коду $\tau^{\prime}$ и множеству $V'$.
	\end{enumerate}
	
	В конце работы алгоритма в $V'$ останется две вершины, они составляют последнее ребро дерева.
	
	Докажем, что на выходе из алгоритма мы получаем граф $G = (V, E)$, являющийся деревом.
	
	У полученного графа $n-1$ вершина. Докажем, что он ациклический, откуда получим по теореме \ref{equivalences}, что граф является деревом.
	
	Предположим противное, т.\:е. в графе есть цикл $Z$. Возьмем первое ребро цикла $Z$, полученное по нашему алгоритму. Пусть это ребро $\left(v_i, v_j \right)$. Рассмотри шаг алгорима, на котором оно было получено.
	Н.Н.О. вершина $v_i$ была взята из множества $V'$. Но тогда на последующих шагах алгоритма не было построено ребер с $v_i$. Действительно, по построению $v_i$ нету ни в оставшейся части кода Прюфера, ни в множестве $V'\setminus{\left\lbrace v_i\right\rbrace }$. Поскольку мы взяли первое ребро в цикле, встретившееся в ходе алгоритма, то ребро с вершиной $v_i$, участвующее в цикле, не могло встретиться и до рассмотренного шага. Противоречие, т.\:е. цикла получиться не могло.
	
	Следовательно, отображение $\varphi$ является биекцией, что завершает доказательство теоремы.
\end{proof}

\bigskip

В качестве бонуса докажем следующее утверждение:

\begin{prop}
Вершина $v_i$ встречается в коде Прюфера $deg(v_i) - 1$ раз.
\end{prop}

\begin{proof}
Докажем, что если в коде Прюфера число $i$ встречается $k$ раз, то в соответствующем ему графе степень $deg(v_i) = k + 1$. Заметим, что в начале на доске выписано $n-2 + n = 2n - 2$ чисел. Выписанные числа - это числа кода Прюфера, а также элементы множества $\left\lbrace 1… n\right\rbrace$, поэтому число $i$ встречается среди выписанных $k+1$ раз. Поскольку по алгоритму код отображается в дерево, среди выписанных ребер нет петель и кратных ребер. Отсюда, результатом алгоритма является разбиение выписанных чисел на пары, т.\:е. $n - 1$ элемент множества $\left(\begin{array}{c} V \\ 2\end{array}\right)$. Отсюда, среди $n-1$ пары ровно $k+1$ тех, в которой участвует вершина $v_i$. 
\end{proof}

    \subsubsection{Унициклические графы}

\begin{Def}
	Граф называется \emph{унициклическим}, если он связен и содержит ровно один цикл.
\end{Def}

\begin{lemma} 
	Число лесов на n вершинах с $r$ деревьями, где i-му дереву ппринадлежит вершина с номером i, равно 
	$ \mathlarger n^{n-1-r}r$
\end{lemma}\label{forestlemma}


Обозначим за $U(n)$ количество унициклических графов на $n$ вершинах. Докажем, что справедлива формула:

$$
U(n)=\sum_{r=3}^{n} C_{n}^{r} \cdot \frac{(r-1)!}{2} \cdot n^{n-1-r} \cdot r
$$

Неформально можно рассуждать так: выберем произвольное число $r$ ($r \geqslant 3$) и посчитаем количество графов с единственным циклом длины $r$. Мы можем выбрать $C_{n}^{r}$ способами $r$ вершин, участвующих в цикле и соединить их в цикле $\frac{(r-1) !}{2}$ способами. Если выкинуть ребра цикла, то останется лес из $r$
деревьев на $n$ вершинах, где $i$ -ое дерево содержит $v_{i} .$ Таких деревьев по \ref{forestlemma} ровно $n^{n-1-r} r$.

Найдем асимптотику для функции $U(n)$.
\begin{theorem}
	$$U(n) \sim \sqrt{\frac{\pi}{8}} n^{n-\frac{1}{2}}$$
\end{theorem} \label{izvrat}

\begin{proof}
	$$
	\begin{aligned}
	U(n) &=\sum_{r=3}^{n} C_{n}^{r} \frac{(r-1) !}{2} n^{n-1-r} r \\
	&=\frac{1}{2} n^{n-1} \sum_{r=3}^{n} n^{-r} \prod_{i = 1}^{r-1}\left(n-i\right) \\
	&=\frac{1}{2} n^{n-1} \sum_{r=3}^{n} \prod_{i = 1}^{r-1}\left(1-\frac{i}{n}\right)
	\end{aligned}
	$$
	Рассмотрим отдельно $\sum_{r=3}^{n}\left(1-\frac{1}{n}\right) \ldots\left(1-\frac{r-1}{n}\right)$. По формуле, использовавшейся в процессе доказательства \ref{exponent}, имеем:
	
	$$
	\sum_{i = 1}^{n} \ln\left(1 - \frac{i}{n}\right)  = -\frac{r(r-1)}{2n} + O\left( \frac{r^3}{n^2}\right) 
	$$
	
	\qquad
	
	Заметим, что при $r>\sqrt{n}$ дробь $\mathlarger{-\frac{r(r-1)}{2 n} \rightarrow-\infty} \Rightarrow e^{\frac{-r(r-1)}{2 n}} \rightarrow 0
	$, при $r<n^{\frac{2}{3}}: O\left(\frac{r^{3}}{n^{2}}\right)=o(1) .$ Тогда имеем
	
	
	\begin{equation*}
	\mathlarger{
	\sum_{r=3}^{n}\prod_{i = 1}^{r-1}\left(1-\frac{i}{n}\right)=} \\
 	\underbrace{\mathlarger{\sum_{r=3}^{\left[n^{0.6}\right]}\prod_{i = 1}^{r-1}\left(1-\frac{i}{n}\right)}}_{S_1}+\underbrace{\mathlarger{\sum_{r=\left[n^{0.6}\right]+1}^{n}\prod_{i = 1}^{r-1}\left(1-\frac{i}{n}\right)}}_{S_2} \footnote{Константа 0.6 взята, поскольу она близка к числу 2/3.}
	\end{equation*}

	Оценим $S_{2}:$
	$$
	\begin{aligned}
		S_{2} \leqslant & \sum_{r=\left[n^{0.6}\right]+1}^{n} e^{-\frac{r(r-1)}{2 n}}  \\
		\leqslant & \sum_{r=\left[n^{0.6}\right]+1}^{n} e^{-\frac{n^{0,6}\left(n^{0,6}-1\right)}{2 n}} \\
		=& \sum_{r=\left[n^{0.6}\right]+1}^{n} e^{-\frac{n^{0,2}(1+o(1))}{2}} < \\
		<& n e^{-\frac{n^{0,2}(1+o(1))}{2}} \rightarrow 0 
	\end{aligned} \label{S_2}
	$$
	
	Теперь оценим $S_{1}$:
	
	\begin{equation*}
	\begin{aligned}
	S_{1} & \sim \sum_{r=3}^{\left[n^{0.6}\right]} e^{-\frac{r(r-1)}{2 n}+O\left(\frac{r^{3}}{n^{2}}\right)} \\
	& \sim \sum_{r=3}^{\left[n^{0.6}\right]} e^{-\frac{r(r-1)}{2 n}} \\
	& \sim \sum_{r=3}^{\left[n^{0.6}\right]} e^{-\frac{r^{2}}{2 n}} \\
	&=\sum_{r=0}^{\infty} e^{-\frac{r^{2}}{2 n}}-\sum_{r=0}^{2} e^{-\frac{r^{2}}{2 n}}-\sum_{r=\left[n^{0.6}\right]+1}^{\infty} e^{-\frac{r^{2}}{2 n}}
	\end{aligned}
	\end{equation*} \label{S2}
	
	$\mathlarger \sum_{r=0}^{2} e^{-\frac{r^{2}}{2 n}} \rightarrow 3$ при $n \rightarrow \infty,$ а $\mathlarger \sum_{r=0}^{\infty} e^{-\frac{r^{2}}{2 n}} \sim \mathlarger \int_{r=0}^{\infty} e^{-\frac{r^{2}}{2 n}} dr$. А это известный всем нам интеграл:
	
	$$
	\begin{array}{c}
	\mathlarger{
	\int_{r=0}^{\infty} e^{-\frac{r^{2}}{2 n}} dr} =\\
	=\mathlarger{\sqrt{n} \int_{x=0}^{\infty} e^{-\frac{z^{2}}{2}} dx} \\
	=\sqrt{n} \frac{\sqrt{2\pi}}{2}=\sqrt{\frac{\pi n}{2}} \\
	\end{array} \label{integral}
	$$
	
	Осталось оценить последнюю из сумм.
	
	$$
	\sum_{r=\left[n^{0.6}\right]+1} e^{-\frac{r^{2}}{2 n}}=\underbrace{\sum_{r=\left[n^{0.6}\right]+1}^{n^{2}} e^{-\frac{r^{2}}{2 n}}}_{S_1^{\prime}} + \underbrace{\sum_{n^{2}+1}^{\infty} e^{-\frac{r^{2}}{2 n}}}_{S_{2}^{\prime}}
	$$
	
	Оценим $S_1^{\prime}$. Аналогично  \ref{S2}, $\mathlarger{S_{1}^{\prime} \leqslant \sum_{r=\left[n^{0.6}\right]+1}^{n^{2}} e^{-\frac{n^{1,2}}{2 n}}<n^{2} e^{-\frac{n^{0,2}}{2}} \rightarrow 0}$. Теперь оценим  $S_{2}^{\prime}$. Заметим, что
	отношение двух соседних слагаемых в $S_{2}^{\prime}$ не превосходит $e^{-n}$. Действительно,
	
	$$ \mathlarger{
	e^{-\frac{(r+1)^{2}-r^{2}}{2 n}}=e^{-\frac{2 r-1}{2 n}}<e^{-\frac{r}{n}}<e^{-n}}
	$$
	
	В таком случае можем оценить $S_{2}^{\prime}$ суммой геометрической прогрессии со знаменателем $e^{-n}$ и первым членом, равным  $e^{-\frac{(n^{2}+1)^{2}}{2 n}}$:
	
	$$
	\mathlarger{ S_{2}^{\prime} < e^{-\frac{(n^{2}+1)^{2}}{2 n}}\left(1+e^{-n}+e^{-2 n}+\ldots\right)= e^{ -\frac{\left(n^{2}+1\right)^{2}}{2 n}} \frac{1}{1-e^{-n}} \rightarrow 0}
	$$
	
	Итого:
	$$
	U(n) \sim \frac{1}{2} n^{n-1}\left(S_{1}+S_{2}\right) \sim \frac{1}{2} n^{n-1} \sqrt{\frac{\pi n}{2}}=\sqrt{\frac{\pi}{8}} n^{n-\frac{1}{2}},
	$$
\end{proof}

Основной результат главы доказан. Посмотрим тепрь на асимптотику еще одной величины, но оставим ее без доказательства.

Обозначим за $C(n, k)$ количество связных графов на $n$ вершинах с k ребрами. Из очевидных соображений имеем $C(n, k) = 0$ для $k = 1, \dots n-2$. Из теоремы \ref{Caley} знаем, что $C(n, n-1) = n^{n-2}$. Из только что доказанной теоремы \ref{izvrat} $C(n, n) = U(n)$. 

\begin{theorem}[б/д]
	$$C(n, n+1) \sim \frac{5}{24}n^{n+1}$$
\end{theorem}


\subsubsection{Эйлеровы и Гамильтоновы графы}

\begin{Def}
	\emph{Эйлеров} цикл (путь) - это несамопересекающийся цикл (путь), проходящий через каждое ребро графа ровно по одному разу.
\end{Def}

\begin{Def}
	\emph{Гамильтонов} путь (цикл) - это несамопересекающийся путь (цикл), проходящий через каждую вершину (кроме начальной и последней) графа ровно по одному разу.
\end{Def}

\begin{Def}
	Граф называется \emph{гамильтоновым}, если существует несамопересекающийся цикл, проходящий по каждой (кроме начальной и последней) вершине ровно по одному разу.
\end{Def}

Есть способы достаточно быстро проверить, вяляется ли граф эйлеровым. А именно, верна следующая теорема:

\begin{theorem}
	Для связного графа следующие утверждения эквивалентны:
	\begin{enumerate}
		\item Граф эйлеров
		\item Степень каждой вершины четна
		\item Множество ребер графа распадается в объединение непересекающихся по ребрам простых циклов
	\end{enumerate}
\end{theorem}

\begin{proof}
	$1 \Rightarrow 2$.
	Если у какой-то вершины степень нечетна, то существует минимум две вершины $v_1$ и $v_2$ с нечетной степенью. Раз так, то наш эйлеров путь начался в $v_1$ и закончился в $v_2$ или наоборот, но он является циклом.
	 
	$2 \Rightarrow 3$.
	Будем идти по графу из произвольной вершины $v$. Поскольку степень каждой вершины четна, мы, войдя в очередную вершину, сможем выйти из нее. будем продолжать движение, пока не вернемся в одну из уже выбранных вершин. Тогда мы найдем некоторый простой цикл.
	Удалим все его ребра из $G$ и получим новый граф, возможно, с несколькими компонентами связности.
	Проделаем аналогичную операцию в остальных компонентах и заметим при этом, что количество ребер ф графе уменьшается, а свойство четности степеней вершин сохраняется.
	
	$3 \Rightarrow 1$.
	Доказательство по индукции. Для одного простого цикла утверждение очевидно.
	Возьмем произвольный граф и удалим один простой цикл $Z$. Полученный граф $G^{\prime}$ распадется на компоненты связности, в каждой из который есть эйлеров цикл по предположению индукции. Начнем обходить граф по исходному циклу $Z$. Попадая на очередном шаге в вершину в одной из компонент связности, обходим все ребра компоненты связности, возввращаясь в иходную вершину цикла $Z$. Таким образом мы обойдем весь граф.
\end{proof}

Задача проверки Гамильтоновости графа является $NP$-полной, т.е. неизвестен алгоритм, который бы в общем случае за полиномиальное время смог проверить, является ли произвольный граф гамильтоновым.
Тем не менее, в частных случаях есть множество способов понять, является ли граф гамильтоновым. 
Сформулируем два признака.

\begin{theorem}(Признак Дирака)
	Пусть $G = (V, E), |V| = n$. Пусть $\forall v \in V deg(v) \geqslant \frac{n}{2}$. Тогда граф Гамильтонов.
\end{theorem}

\begin{Def}
	Подмножество $W \subseteq G$ называется \emph{независимым}, если $\forall x, y \in W : (x, y) \notin E$. 
\end{Def}

\begin{Def}
	$\alpha(G) = \max \left\lbrace k \in \mathbb{N} \mid \exists W \subset G : \forall x, y \in W  (x, y) \notin E \land |W| = k  \right\rbrace$ - \emph{число независимости} графа. В частности, $\omega(G)$ обозначается одно из максимальных независимых множеств графа.
\end{Def}

\begin{Def}
	$\varkappa(G) = \min \left\lbrace k \in \mathbb{N} \mid \exists W \subset G : |W| = k \land G|_{V\setminus E} - \text{несвязен} \right\rbrace$ - \emph{вершинная связность} графа. 
\end{Def}

\begin{theorem}[Эрдеш, Хватал]
	Пусть $G(V, E)$ - граф и $\alpha(G) \leqslant \kappa(G)$. Тогда $G$ гамильтонов.
\end{theorem}

    \begin{theorem}[Эрдеш, Хватал]
	Пусть $G(V, E)$ - граф, $ |V| \geqslant 3 $ и $\alpha(G) \leqslant \varkappa(G)$. Тогда $G$ -- гамильтонов.
\end{theorem}

\begin{proof}
	\:
	\begin{enumerate} \renewcommand{\theenumi}{\bfseries\arabic{enumi}}
		\item Предположим, что в $ G $ нет циклов. Поскольку $1 \leqslant \alpha(G)$ (множество из одной вершины всегда независимое), то $ \alpha(G) \leqslant \varkappa(G) \Rightarrow 1 \leqslant \alpha(G)$, следовательно, мы должны удалить хотя бы одну вершину, чтобы $ G $ стал несвязен. Отсюда $ G $ связен и без циклов, т.е. дерево. Для $ |V| < 3$ утверждение очевидно. Если же $ |V| \geqslant 3$, то в дереве найдется две висячих вершины. Две вершины являются независимым множеством мощности 2, откуда $ \alpha(G) \geqslant 2$. Но в дереве на $ \geqslant 3$ вершинах, очевидно, есть и вершина степени $ \geqslant 2$. Её удаление разбивает граф на компоненты связности, откуда $ \varkappa = 1$. Получаем противоречие.
		
		\item Пусть в $ G $ есть цикл. Рассмотрим самый длинный простой цикл и предположим, что он не покрывает все вершины: $ C = {x_1, \ldots, x_k}, k < |V| = n $. Удалим вершины цикла и обозначим новый граф за $ G^\prime $. Пусть $ W $ -- множество вершин одной из компонент в $ G^\prime $. Обозначим $ N_{w}\left(G \right) = \left\lbrace y \in V\setminus W : \exists z \in W : (y, z) \in E(G)\right\rbrace$. т.\:е. множество соседей вершин из $ W $ в исходном графе $ G $.
		
		%\begin{wrapfigure}{r}{0.2\textwidth}
		%	\centering
		%	\begin{tikzpicture}
		%		
		%		\def \n {8}
		%		\def \k {6}
		%		\def \radius {2.5cm}
		%		\def \margin {10} % margin in angles, depends on the radius
				
		%		\foreach \s in {1,...,\n}
		%		{
		%			\node[draw, circle] at ({360/\n * (\s - 1)}:\radius) %{$x_{\s}$};
		%			\draw[->, >=latex] ({360/\n * (\s - 1)+\margin}:\radius) 
		%			arc ({360/\n * (\s - 1)+\margin}:{360/\n * %(\s)-\margin}:\radius);
		%		}
		%	\end{tikzpicture}
		%\end{wrapfigure}
	
		\begin{prop}
			\:
			\begin{enumerate} \renewcommand{\theenumi}{\arabic{enumi}}
				\item $ N_{w}\left(G \right) \subseteq C$
				\item $ N_{w}\left(G \right) \neq C$
				\item $ \varkappa(G) \leqslant |N_w(G)| $
			\end{enumerate} \label{prop:N_w}
			
		\end{prop}
		
		\begin{proof}
			\:
			\begin{enumerate} \renewcommand{\theenumi}{\arabic{enumi}}
				\item Если предположить, что есть вершина не из $ C $, она после удаления $ С $ осталась в $ W $, противоречие.
				\item Если нашлись две вершины $ x_i, x_{i + 1} \in N_w(G) $, то получим более длинный цикл, заменив в $ C $ кусок $\dots \rightarrow x_i \rightarrow x_{i + 1} \rightarrow \dots$ на $ \dots \rightarrow x_i \rightarrow c \rightarrow x_{i + 1} \rightarrow \dots$.
				\item Когда мы удаляем $ N_w(G), $ по пункту (b) остается хотя бы одна вершина цикла $ C $, не соединенная с множеством $ W $, а также остается сам $ W $, т.е. граф распадается на компоненты связности.
			\end{enumerate}
		\end{proof}
		Обозначим $ M = \left\lbrace x_{i + 1} : x_i \in N_w(G)\right\rbrace $, другими словами, это - множество вершин цикла $ C $, родители которых соединены с $ W $. Очевидно, $ |M| = |N_w(G)| $
		\begin{prop}
			$ M $ - независимое множество.
		\end{prop}
		\begin{proof}
			Предположим, что множество $ M $ не независимое, т.е. $ \exists\: x_{i + 1}, x_{j + 1} \in M : (x_{i + 1}, x_{j + 1}) \in E $ и $ x_i, x_j \in N_w(G) $. Пусть вершины $ x_{i}, x_{j} $ соединены с какими-то вершинами $ a, b $ из $ M $ соответственно ($ a $ и $ b $ могут совпадать).
			
			
			Имеем исходный цикл $ C $ : $x_1 \rightarrow \dots \rightarrow x_i \rightarrow x_{i + 1} \rightarrow \dots \rightarrow x_j \rightarrow x_{j + 1} \rightarrow \dots \rightarrow x_n$. Можем рассмотреть новый цикл $ C^\prime $ :
			$x_1 \rightarrow \dots \rightarrow x_i \rightarrow a \rightarrow \dots \rightarrow b \rightarrow x_j \rightarrow \ldots \rightarrow x_{i + 1} \rightarrow x_{j + 1} \dots \rightarrow x_n \rightarrow x_1$. Он более длинный, поскольку добавились как минимум три ребра, а из исходного цикла были взяты все ребра, кроме $ x_ix_{i + 1}, x_jx_{j + 1} $. Это противоречит выбору $ C $ как наибольшего. Следовательно, $ M $ действительно независимое.
			
		\end{proof}
		Получаем, что $ \alpha(G) \geqslant |M|$. Осталось заметить, что по построению $ M $ - точки, идущие в цикле $ C $ в выбранном порядке нумерации сразу за точками $ N_w(G) $. Две соседние точки цикла по пункту ($ b $) утверждения \ref{prop:N_w} не могут вместе лежать в цикле. Следовательно, вершины из $ M $ не соединены с $ W $. Поскольку мы изначально предположили, что $ W $ непусто, то, добавив к $ M $ произвольную точку $ w \in W$, получим $ M' = M \cup \left\lbrace w\right\rbrace $ - независимое множество. Отсюда $ \alpha(G) \geqslant  |N_w(G)| + 1 $. По пункту ($ c $) утверждения \ref{prop:N_w} $ \varkappa(G) \leqslant |N_w(G)| $, откуда $ \alpha(G) > \varkappa(G) $. Противоречие.
			
	\end{enumerate}
\end{proof}

\begin{Def}
	Граф называется \emph{регулярным}, если степени всех его вершин равны.
\end{Def}

Рассмотрим следующий интересный регулярный граф и докажем, что он удовлетворяет признаку Эрдеша-Хватала.

Обохначим $ N = \left\lbrace 1, 2, \ldots n\right\rbrace $.Рассмотрим $ G = (V, E) : V = \left\lbrace A \subset N  : |A| = 3 \right\rbrace, E = \left\lbrace (A, B) : |A\cap B| = 1\right\rbrace  $. Другими словами, $ G $ -- граф на множестве трехэлементных подмножеств $ N $, где два множества соединены ребром, если они пересекаются по одному элементу.

Очевидно, $ |V| = C_n^3 $. Найдем $|E|$. Степень каждой вершины $  G $ -- $3 \cdot C_{n - 3}^2$, поскольку для каждого из трех элементов $ {a, b, c} $ мы можем выбрать двухэлементное подмножество $ Q \subset N \setminus \left\lbrace a, b, c \right\rbrace $ и объединить их. Очевидно, таким образом описываются все множества, смежные с вершиной  $ \left\lbrace a, b, c\right\rbrace $ в $ G $.

Поскольку $  G $ регулярен,  $ |E| = \frac{1}{2} \cdot |V| \cdot 3 \cdot  C_{n - 3}^2  = \frac{1}{2} C_n^3 \cdot 3 \cdot  C_{n - 3}^2 $. Асимптотически

$$ |E| \sim \dfrac{3}{2} \cdot \dfrac{n^3}{6} \cdot \dfrac{n^2}{2} =\dfrac{n^5}{8} $$.

Вспомним, что признаку Дирака удволетворяют графы, у которых $ m $ вершин и не меньше, чем $ \dfrac{m^2}{4} $ ребер.
Поскольку в нашем случае $|V| = O(n^3)$, а $|E| = O((n^3)^{\frac{5}{3}}) $, при больших $ n $ наш граф не будет удовлетворять признаку Дирака.

Убедимся, что $ G $ удовлетворяет признаку Эрдеша-Хватала. 
Посмотрим на $ \alpha(G) $. Укажем $ W \subset V : W = {A_1, \ldots A_k} : |A_i \cap A_j| \neq 1 $.
Каждому множеству $ A_i $ сопоставим вектор $ \overrightarrow{x_i} $ в $ \left\lbrace 0, 1\right\rbrace ^{|V|} $, где на позиции $ i $ стоит 1, когда $ i \in W $, и 0 иначе.
Отметим, что $ |A_i \cap A_j| = (\overrightarrow{x_i}, \overrightarrow{x_j}) $.

\begin{lemma}
	$ \overrightarrow{x_1}, \ldots \overrightarrow{x_k} $ линейно независимы над полем $ \Z_2 $.
\end{lemma}  

\begin{proof}
	Предположим, что $ c_1 \overrightarrow{x_1} + \ldots + c_k \overrightarrow{x_k} = 0 $, где $ c_i \in \Z_2 $.
	Скалярно умножим выражение на $ \overrightarrow{x_i} $ для $ i = 1, \ldots k $.
	Для $ i = 1 $ получим $ c_1 (\overrightarrow{x_1}, \overrightarrow{x_1}) + \ldots + c_k (\overrightarrow{x_k}, \overrightarrow{x_1}) = 0$. Заметим теперь, что первое слагаемое равно $ 3 c_1 $, т.е. $ c_1 $ в поле $ \Z_2 $. Все оставшиеся слагаемые равны 0 или 2, т.е. 0 в нашем поле. Действительно, $ (\overrightarrow{x_j}, \overrightarrow{x_1}) \neq 1 $, поскольку векторы соответствуют множествам, не пересекающимся по 1 элементу. $ (\overrightarrow{x_j}, \overrightarrow{x_1}) \neq 3 $, поскольку $ j \neq 1 $. Следовательно, $ (\overrightarrow{x_j}, \overrightarrow{x_1}) = 0 $ над $ \Z_2 $. Получаем равенство $ c_1 = 0$, откуда $ c_1 = 0 $. Умножая скалярно выражение на $ x_i $ для остальных $ i $, получаем, что $ \forall i \; c_i = 0 $.
\end{proof}

\begin{corollary}
	$ \alpha(G) \leqslant |V| $.
\end{corollary} \label{cor:alpha}

\begin{proof}
	Количество линейно независимых векторов не превосходит размерности линейного пространства. 
\end{proof}

Несложно привести конструкцию для $ n - 2 $ независимых множеств.

\begin{center}
	\begin{tabular}[h!]{cccccc} 
		1 & 2 & 3 &&& \\ 
		1 & 2 & & 4 && \\ 
		1 & 2 & & & 5 & \\
		& $\ldots$ & & $ \ldots $ &&  \\
	\end{tabular}
\end{center}

Эту оценку можно в некоторых случаях улучшить.

\begin{theorem}
	$$
	\alpha(G)=\begin{cases}
		n, & \; n \equiv 0\bmod 4\\
		n - 1, & \; n \equiv 1 \bmod{4}\\
		n - 2, & \; \text{иначе} \\
	\end{cases}
	$$
\end{theorem}

\begin{proof}
	Пример для каждого случая приводится конструкцией разбиения множества на 4-элементные множества, возожно, с остатком меньшей мощности. В каждом из них можно взять по 4 трехэлементых множества, каждые два из которых будут пересекаться по двум элементам.
	\begin{enumerate} \renewcommand{\theenumi}{\arabic{enumi}}
		\item Если $ n = 4k $, получаем в точности $ n $ независимых множеств.
		\item Если $ n = 4k + 1 $, получаем $ 4k = n-1 $ независимых множеств.
		\item Если $ n = 4k + 2 $, получаем $ 4k = n-2 $ множеств. Наконец, для $ n = 4k + 3 $ после разбиения получится 3-элементное множество в остатке, которое мы также возьмем в набор. Т.е. получим $ 4k + 1 = n-2 $ множеств.
	\end{enumerate}

Осталось показать, что в случаях 2 и 3 нельзя получить лучшей оценки.
\end{proof}

Вернемся к доказательству оценки $\alpha(G) \leqslant \varkappa(G)$. По \ref{cor:alpha} имеем $\alpha(G) \leqslant n $. Осталось показать, что $n \leqslant \varkappa(G) $.

\begin{lemma}
	Пусть дан $ G = (V, E) $. Пусть $ u, v \in V $. Обозначим за $ f(u, v) $ количество общих соседей для $ u $, $ v $. Тогда $ \varkappa(G) \geqslant \mathlarger {\min_{v, u \in V} f(u, v) }$
\end{lemma} \label{lemma:neigh}

 \begin{proof}
 	Для удобства $ \mathlarger \min_{v, u \in V} f(u, v) $ обозначим за $ k $.
 	Если удалить из графа не более $ k - 1 $ вершин, то для любых двух вершин $ u, v $ количество общих соседей не меньше $ k $, а следовательно между ними есть путь.
 \end{proof}



В следующей табличке приводятся случаи количества общих соседей у вершин нашего графа.

\begin{center}
	\begin{tabular}[h!]{|ccc|}
		\hline
		
		Тип множеств по мощности пересечения & Кол-во $ N $ общих соседей & Асимптотика $ N $ \\
		\hline 
		$  \left\lbrace 1, 2, 3 \right\rbrace, \left\lbrace 4, 5, 6 \right\rbrace, \; |\cap| = 0 $ & $ 9(n - 6)$ & $ 9n $ \\ 
		$  \left\lbrace 1, 2, 3 \right\rbrace, \left\lbrace 3, 4, 5 \right\rbrace, \; |\cap| = 1 $ & $ C_{n-5}^2 + 4(n-5) $ & $ n^2 $ \\
		$  \left\lbrace 1, 2, 3 \right\rbrace, \left\lbrace 2, 3, 4 \right\rbrace, \; |\cap| = 2 $ & $ 2C_{n-4}^2 + n - 4 $ & $ n^2 $  \\
		\hline
	\end{tabular}
\end{center}

\begin{note}
	Случай $ |\cap| = 3 $ мы не рассматриваем, поскольку количество общих с собой соседей у произвольной вершины равно ее степени и заведомо не меньше, чем количество общих сосдей с любой другой вершиной.
\end{note}

Следовательно, по лемме \ref{lemma:neigh} при больших $ n $  $ \varkappa(G) \geqslant 9n $, т.е. примерно в 9 раз больше $ \alpha(G) $. Следовательно, начиная с какого-то $ n $ граф гамильтонов. 

\subsubsection{Граф  $G(n, r, s)$}

Читатель уже мог догадаться, что рассмотрненный в предыдущем параграфе граф суть частный случай графа, строящегося на подмножествах $n$-элементного множества одной мощности. Дадим соответсвующие определения.

\begin{Def}
	Обозначим $ N = \{1, \ldots, n\}$. Рассмторим граф $ G = (V, E) : V = \left\lbrace A \subset N  : |A| = r \right\rbrace, E = \left\lbrace (A, B) : |A\cap B| = s\right\rbrace $. Такой граф обозначается $ G(n, r, s) $. 
\end{Def}

\begin{prop}
	 B графе $ \mathlarger {G(n, r, s) : |E|=\frac{1}{2} C_{n}^{r} C_{r}^{s} C_{n-r}^{r-s} }$
\end{prop}

\begin{proof}
	Поскольку граф регулярный, то количество вершин считается по формуле $ \frac{1}{2} |V| \cdot \deg(v)$. $|V| = C_{n}^{r}$ - кол-во $r$-элементных подмножеств $ N $. $C_{r}^{s}-$ кол-во способов
	выбрать $s$ элементов из произвольного $ r $-элементного множества, а $C_{n-r}^{r-s}-$ кол-во способов добрать оставшиеся элементы во множество.
\end{proof}

%Утверждение 1.1.1.2. Количество треугольников в графе $G(n, r, s)$ равно
%$$
%\frac{|E|}{3}\left(\sum_{i=0}^{s} C_{s}^{i} C_{r-s}^{s-i} C_{r-s}^{s-i} C_{n-2 r+s}^{r-2 s+i}\right)
%$$
%$ G(n, r, s) $ 

\begin{brainer}
	Для $ n = 4k $ рассмотрим $ G(n, \frac{n}{2}, \frac{n}{4}), n = 4k $.
	\begin{enumerate} \renewcommand{\theenumi}{\arabic{enumi}}
		\item Найдите число треугольников в этом графе.
		\item Какой объект соответствует максимальной клике в этом графе?
	\end{enumerate}
\end{brainer}









    \subsubsection{Последовательности де Брёйна}

\begin{Def}
	\emph{Последовательностью де Брёйна} для числа $ n $ называется последовательность $ a_1, a_2, \ldots a_N, \;a_i \in \{0, 1\} $, так что все среди подслова длины $ n $ встречаются все различные слова длины $ n $ над афлавитом $ \{0, 1\} $.
\end{Def}

Поскольку число слов длины $ n $ над афлавитом $ \{0, 1\} $ равно $ 2^n $, то $ N = 2^n + n - 1 $.

Перейдем к доказательству существования для произвольного $ n $ и построению последовательности де Брейна. Приведем два способа построения.

\begin{enumerate} \renewcommand{\theenumi}{\bfseries \arabic{enumi}}
	\item {\bfseries Граф де Брейна} Первый способ построения заключается в построении графа и нахождения в нем Эйлерова цикла.
	\begin{Def}
		\emph{Графом Де Брёйна} называется ориентированный граф на множестве слов длины $ n-1 $ над алфавитом $ \{0, 1\} $, т.ч. есть ребро из слова $ u $ в слово $ v $, когда существует слово  длины $ n $, такое что его префикс - это $ u $, а суффикс - $ v $.
	\end{Def}

	\begin{prop}
		Рассмотренный граф Эйлеров.
	\end{prop}

	\begin{proof}
		Необходимо показать, что для каждой вершины $u$, соответствующей слову $ b_1\ldots b_{n-1} $, $ indeg(u) = outdeg(u) $, откуда по критерию Эйлеровости получим, что он Эйлеров. Слово $ u $ можно дополнить справа 0 или 1, поэтому из него выходит два ребра (возможно, одно из них идет в себя). По тем же соображениям в него входит 2 ребра. Следовательно, утверждения доказано.
		%$$ \begin{cases}
		%		b_2\ldots b_{n-1} 0 \neq b_1\ldots b_{n-1} \\
		%		b_2\ldots b_{n-1} 1 \neq b_1\ldots b_{n-1}
		%\end{cases}  $$
	\end{proof}

	\begin{prop}
		Произвольный Эйлеров цикл в рассмотренном графе задает правильную последовательность, а именно, если выписать слово, соответствующее стартовой вершине, а на каждом переходе в новую вершину писать его последнюю букву, то получится последовательность де Брёйна.
	\end{prop}

	\begin{proof}
		Слово, соответствующее каждой из вершин графа, встретится в искомой последовательности. Действительно, на первом шаге мы выписали слово, соответствующее начальной вершине, а на очередном выписывании буквы мы получали префикс, равный слову, соответствующему вершине, в которую мы перешли. Поскольку длина полученной последовательности равна $ (n - 1) + 2^n = N$, то никакое слово не встретится дважды.
	\end{proof}

	\begin{note}
		\href{https://mipt.lectoriy.ru/file/synopsis/pdf/Maths-DiscreteMathem-M05-Raygor-131016.03.pdf}{Здесь} вы можете найти наглядный пример построения и обхода графа для $ n = 3 $.
	\end{note}

	\item {\bfseries Правило "0 лучше 1".} Второй алгоритм значительно проще реализуется - но сложнее доказывается.
	Вначале выставим $ n $ единиц.
\end{enumerate}

    
    %\newpage
\section{Additional information}
\end{document}